\section{Finite observations and detection of unseen area}\label{sec:statistics}
\subsection{Density estimation and simultaneous decisions}
The statistical experiment keeps the same marked endpoint sensor. A
preparation at one window supplies one endpoint cell or a failure/outside
outcome. Its indicators for all cells are obtained from the same trial;
there is no separate preparation for each cell.

\begin{lemma}[Two derivatives from cell masses]\label{lem:histogram}
On nested fixed rectangles, a $C^{2,\beta}$ density with norm at most
$M$ has a linear reconstruction from square-cell probabilities of side
$h$ satisfying
\begin{equation}\label{eq:histogram-error}
 \|\widehat f-f\|_{C^2}\le C(h^\beta+\varepsilon h^{-4})
\end{equation}
when every cell probability has error at most $\varepsilon$. There
are $O(h^{-2})$ cells. With $N$ independent trials at each window in
a list of $J$ records, the simultaneous bound is
\begin{equation}\label{eq:bernstein-density}
 e_N\le C\{h^\beta+\sqrt{r_N}h^{-3}+r_Nh^{-4}\},\qquad
 r_N=\frac{\log(CJN/\delta)}{N},
\end{equation}
at probability at least $1-\delta$, for the mesh chosen below.
\end{lemma}
\begin{proof}
The averages of $1,z,z^2$ on the three unit cells centered at $-1,0,1$
are the rows $(1,i,i^2+1/12)$ of an invertible matrix. Its tensor square
recovers a coordinatewise quadratic polynomial from a $3\times3$ cell
block. The degree-two Taylor polynomial of a $C^{2,\beta}$ function
has remainder $O(h^{2+\beta})$ on such a block. Polynomial reproduction
and rescaling give derivative errors $O(h^{2+\beta-j})$, $j\le2$.
A probability error $\varepsilon$ is an average-density error
$\varepsilon h^{-2}$ and becomes $O(\varepsilon h^{-2-j})$ after
$j$ derivatives. Patch the local polynomials with a smooth partition of
unity of bounded overlap, whose derivatives of order $j$ are $O(h^{-j})$.
Leibniz's rule preserves these estimates. Use a larger observation
rectangle to provide all boundary stencils. This proves
\eqref{eq:histogram-error}.

Each cell probability is at most $Mh^2$. Bernstein's inequality for
its Bernoulli indicator bounds every cell error by
$C(\sqrt{h^2r_N}+r_N)$ after a union bound over both windows and
all records. The mesh $h=r_N^{1/(2\beta+6)}$ has $h^{-2}\le N$
for sufficiently large $N$, so $\log(CJN/\delta)$ dominates the
number of cell events in this union bound. Substitution proves
\eqref{eq:bernstein-density}. Dependence between different cell
indicators from a single trial is irrelevant to a union bound.
\end{proof}

A fixed pilot suffices to place the active windows. Under the local
priors there is a common right collar $d_*>0$ with
$cd^2\le P(T_0+d)\le Cd^2$ and zero mass before $T_0$. Choose a
small fixed threshold $\lambda>0$ and a fixed grid in each known onset
bracket, extended a short distance beyond it. Take
$\sqrt{4\lambda/c}<d_*/32$, grid spacing less than $d_*/32$, and
trigger at a reported mass exceeding $2\lambda$. With pilot errors
below $\lambda/4$, the first trigger $\tau$ lies in
$(T_0,T_0+d_*/16)$. Use $T_1=\tau+d_*/4$ and
$T_2=\tau+d_*/2$ and a rectangle on which $W-T_0\le d_*/8$.
Then both windows lie in the local collar and are strictly active with
a uniform margin. This argument needs no global monotonicity of the
three-impact event. It neither supplies nor estimates an exact onset.
The number of pilot queries is a prior-dependent constant per record.

\begin{proof}[Proof of the finite-histogram assertion of
Theorem~\ref{thm:stat-main}]
Use fresh trials at the two windows after the pilot. Conditional on
the pilot, their times are fixed and the same concentration bound
applies. The pilot means have an ordinary bounded-variable concentration
bound, with their fixed number absorbed in $C$. Thus $CJN$ preparations
suffice for the simultaneous event. The local physical fitting described
in Section~\ref{sec:uniform} changes $e_N$ by at most a constant.

Take $h=r_N^{1/(2\beta+6)}$. The bias and square-root variance terms
in \eqref{eq:bernstein-density} are both
$r_N^{\beta/(2\beta+6)}$; the linear term has exponent
$(2\beta+2)/(2\beta+6)$ and is smaller when $r_N<1$.
Theorem~\ref{thm:uniform-main} applies once this density error is below
its deterministic threshold and gives \eqref{eq:rate-main}. On the
same event the yes/no completeness decision is correct for every
connected component, even when the list contains no covering witness.
Any data-dependent deletion from this finite measured list is allowed
because all record errors were controlled simultaneously. This is not
a claim about a law selected after inspecting unboundedly many uncontrolled
records.
\end{proof}

\begin{corollary}[Repeated inspection and random loss]\label{cor:sequential}
At inspection epoch $m$, use fresh measurements and error allocation
$\delta_m=6\delta/(\pi^2m^2)$. The probability of any false completeness
certificate over all epochs is at most $\delta$, provided all inspected
records are genuine and satisfy the uniform model. No actual witness
retention is assumed for this soundness statement.

Suppose, additionally, an acquisition stream has a fixed finite saturated
covering set $\mathcal S$ of $s$ record orbits, and each not-yet-retained
member is encountered at the next epoch with conditional probability
at least $p>0$. Retain accumulated records. After $m$ epochs the chance
that $\mathcal S$ has not all appeared is at most $s(1-p)^m$.
Once they have appeared and measurement precision meets the fixed threshold,
the certificate accepts. The estimator does not need their identities.
\end{corollary}
\begin{proof}
Sum the epoch failure probabilities. For any one member of $\mathcal S$,
iterating the conditional failure bound gives probability at most
$(1-p)^m$ of its absence; union bound over $s$ members gives the second
assertion. Within-epoch independence between members is unnecessary.
After their arrival the accumulated component covers and is saturated,
so the deterministic certificate accepts on the accuracy event.
\end{proof}

The encounter condition is an explicit assumption about the producer of
records, not a theorem that unmarked trajectories discover normal channels.
Likewise, the probability bound counts phase preparations for measuring
records, not the geometric cost of producing them. The improvement is
that retention need not be promised for a particular observed list:
sufficiency of that list is checked by the area-period certificate.

\subsection{A sharp physical missing-area subexperiment}
The lower area bound $a_0$ is not a harmless normalization. The following
experiment identifies the decision scale when small unseen bodies are
allowed. It also specifies precisely the scope of the statistical
sharpness statement.

\begin{theorem}[Detection of an unseen obstacle]\label{thm:area-testing}
There is an analytic asymmetric periodic table $\mathcal T_0$ and a
family $\mathcal T_a$, $0<a<a_1$, obtained by adding one unobserved
analytic asymmetric obstacle of area $a$ per cell, with the following
properties. A fixed finite collection of selected local return windows
has the same actions and conditional successful endpoint distributions
throughout the family. Its success masses are
\begin{equation}\label{eq:mass-tilt}
                p_j(a)=p_j(0)\frac{A_*}{A_*-a},
\end{equation}
where $A_*$ is the free area of $\mathcal T_0$. Its visible cycle group
is the full fixed period lattice, and its completeness defect is exactly $a$.

For any test of $\mathcal T_0$ against $\mathcal T_a$ using at most
$N$ independent phase preparations, with even adaptive choices among
these finitely many windows, the maximum of its two error probabilities
is at least $1/4$ whenever $a\le cN^{-1/2}$. Conversely, in this
subexperiment the area is estimable to error
$C\sqrt{\log(2/\delta)/N}$ at probability at least $1-\delta$,
and alternatives above a corresponding threshold can be detected.
The constants depend on the fixed experiment. This is not a lower bound
for an arbitrary trajectory sensor or a whole-table minimax theorem.
\end{theorem}
\begin{proof}
Use the one-body square-lattice example in Proposition~\ref{prop:examples},
with a finite set of clear return branches generating the lattice and
small fixed endpoint patches and windows. Their successful phase tubes,
including the initial and final residual-time pieces, can be confined
to arbitrarily thin neighborhoods of finitely many segments. Choose a
point in the complement on the torus and add a sufficiently small scaled
copy of a different asymmetric analytic body there. Scaling realizes
every sufficiently small area $a>0$. It remains disjoint from all the
successful tubes and does not modify any selected branch, action or
twist. The new shape is noncongruent to the original one.

The unnormalized successful phase sets and their endpoint pushforwards
are consequently unchanged, while the free area changes from $A_*$
to $A_*-a$. The local density formula proves \eqref{eq:mass-tilt} and
equality of the conditional successful marks. The outside/failure outcome
is kept as one atom in this experiment. Choose the windows so that all
$p_j(0)$ lie strictly between zero and one and keep $a_1$ small. Every
one-trial relative entropy is therefore exactly that of its Bernoulli
success indicator. Since
\[
 \KL(\operatorname{Bern}(p)\Vert\operatorname{Bern}(q))
                  \le\frac{(p-q)^2}{q(1-q)},
\]
which follows from $\log x\le x-1$, and
$|p_j(a)-p_j(0)|\le Ca$, each one-trial divergence is at most $Ca^2$.
The chain rule gives total divergence at most $CNa^2$ for adaptive
choices from the same finite menu. Pinsker's inequality and the elementary
two-point testing inequality imply maximum testing error at least
$(1-\sqrt{CNa^2/2})/2$. Decreasing $c$ proves the lower bound.

For the upper bound choose one fixed window, with known null success
probability $p_0>0$. Its empirical proportion $\widehat p$ has
Hoeffding error $\sqrt{\log(2/\delta)/(2N)}$. On
$\widehat p\ge p_0/2$, inversion of \eqref{eq:mass-tilt} gives
\[
               \widehat a=A_*\left(1-\frac{p_0}{\widehat p}\right),
\]
clipped to the allowed area interval. This map has a bounded derivative
on the indicated range. The exceptional lower-tail event is covered by
the same concentration bound for sufficiently large $N$. This proves
the estimation and testing assertions. Finally the visible types and
cycle group have not changed, so \eqref{eq:defect} gives defect $a$.
\end{proof}

These alternatives leave the local observed geometry and its derivative
bounds fixed, but the unseen body's area tends to zero. They need not
satisfy the uniform global small-body exclusion in
Section~\ref{sec:uniform}. Thus there is no conflict with that section's
fixed-gap decision. The theorem proves the sharp scale only for this
physical area-testing family; no sharp exponent is inferred for analytic
continuation, shape matching or the count-only germ problem.
